% ------------------------------------------------------------
% §2  History, uses, and the epistemological situation
% ------------------------------------------------------------
\section{The object, observed from outside}\label{sec:history}

Before any internals, here is the object as its users see it. \SHA{} is a
single, fixed, publicly specified function
\[
  \SHA : \bigcup_{0 \,\le\, \ell \,<\, 2^{64}} \bitstrings^{\ell}
         \;\longrightarrow\; \bitstrings^{256},
\]
published by the U.S.\ National Institute of Standards and Technology in
2002 and unchanged since~\citep{FIPS180-4}. It takes a bit string of any
length below $2^{64}$---the empty string, a single bit, the concatenated
works of a library---and returns $256$ bits. (We stress \emph{bit} string:
the standard defines the function on arbitrary bit lengths, not merely on
whole bytes; software interfaces usually round to bytes, but the
mathematical object does not. This matters in \cref{sec:geometry}.)

Two set-theoretic observations frame everything that follows. The domain
is astronomically larger than the codomain, so the function is nowhere
near injective: by pigeonhole, colliding pairs exist in unimaginable
abundance---the average fiber over a $256$-bit output is itself a set of
size around $2^{2^{64}-256}$. Whether the function is \emph{surjective},
on the other hand, is not known: nobody has proved that every $256$-bit
string is attained, and more generally nothing nontrivial is proved about
the image size or about how evenly the domain is spread across the
range---the property Bellare and Kohno isolate as hash-function
\emph{balance}, showing that deviations from it would measurably
accelerate collision search~\citep{BellareKohno2004}; for \SHA{} the
balance is simply unknown. Heuristically, of course, a random function
with this domain misses a given output with probability too small to
name, and \cref{sec:randommaps} is about taking such heuristics seriously
as measurements.

The function is deterministic and involves no secrets: everyone can evaluate
it, and everyone computes the same values.\footnote{``No secrets'' means
no secret \emph{inputs}; it emphatically does not mean no design
intelligence. Every constant and structural parameter was chosen with
visible care, by a designer---the NSA---demonstrably in possession of
nonpublic cryptanalytic knowledge. The clearest episode: NIST published
the family's ancestor, now called SHA-0, in 1993; within two years the
NSA had it withdrawn, citing an undisclosed weakness, and reissued with
exactly one change, a single one-bit rotation added to the message
schedule~\citep{FIPS180-1}. In 1998 academic cryptanalysis independently
discovered a differential collision attack on SHA-0 that this tweak
happens to block~\citep{ChabaudJoux1998}, and by 2005 explicit SHA-0
collisions had been computed~\citep{BihamChen2005}, while the amended
SHA-1 stood for another decade (\cref{sec:falling}). Whatever process
selects these functions, it is not simple random choice; \cref{sec:nums}
discusses how the SHA-2 generation made its constants auditable in
response to exactly this kind of asymmetry.} It is also fast, in a way
worth pausing on: the definition is in effect a small fixed circuit of
bitwise operations and $32$-bit additions, so it is directly expressible
in hardware. A laptop core computes it at gigabytes per second; dedicated
silicon does far better---purpose-built \SHA{} ASIC designs were an
academic subject within two years of
standardization~\citep{Dadda2004}, and the economics of proof-of-work
mining (\cref{sec:uses}) later industrialized them, the arc from CPU
software to modern mining ASICs gaining roughly six orders of magnitude
in throughput per device and five in energy per hash~\citep{Taylor2017}.

What is it \emph{for}? The honest answer is: for behaving, in public, like
a random function---a function sampled uniformly at random from the set of
all functions into $\bitstrings^{256}$---while actually being a short,
fixed program. All of its uses are corollaries of facets of that one
pretension.

\subsection{A catalogue of uses}\label{sec:uses}

\begin{description}
  \item[Fingerprinting.] If $H$ behaves like a random function, two
  different documents have the same $256$-bit hash with probability
  $2^{-256}$; a hash is a practically unforgeable fingerprint. The
  version-control system Git addresses every object it stores by its hash:
  file contents, directory trees, commits. Equality of fingerprints is
  treated as equality of content---mathematically unjustified and, at
  planetary scale, empirically flawless so far.

  \item[Commitment.] Publish $H(x)$ today, reveal $x$ next month; anyone
  can verify you did not change your mind, yet the hash revealed nothing
  about $x$. A random-looking function gives sealed envelopes without
  paper. (This is how one archives a prediction, a bid, or a claim of
  priority.)

  \item[Merkle trees.] Hash the leaves; hash the concatenations of sibling
  hashes; repeat. The root is a $256$-bit commitment to an arbitrarily
  large collection, and membership of any single leaf is certified by a
  logarithmic-length path~\citep{Merkle1988,Merkle1989}. Certificate Transparency
  (the public audit log of all TLS certificates on the web) and every
  blockchain are Merkle trees at heart. Notice the epistemic shape of
  what the tree provides: the inclusion of one transaction among
  billions is \emph{proved} by a chain of a few dozen hashes---a
  certificate anyone can verify in microseconds, revealing essentially
  nothing about the other leaves. Pushed further, this seed grows into
  the theory of zero-knowledge and succinct proof
  systems~\citep{GMR1989}: modern hash-based protocols such as STARKs
  commit to the execution trace of an entire computation with a Merkle
  tree, then prove its correctness by opening a few hash
  chains~\citep{BenSasson2018}. The pattern recurs everywhere in
  mathematical cryptography and deserves explicit statement: a task---%
  here, inverting or colliding \SHA{}---is \emph{presumed} hard, and the
  presumed hardness then serves as a primitive from which richer
  structures (commitments, proofs, signatures, digital money) are built
  by explicit reduction. Hardness conjectures function the way axioms do
  elsewhere in mathematics: unproved, load-bearing, and generative.

  \item[Message authentication.] With a shared secret key $k$, the
  construction $\mathrm{HMAC}(k, m)$, built from two nested \SHA{} calls,
  authenticates messages; its analysis~\citep{BCK1996} is a rare case of a
  \emph{theorem} in this landscape---assuming the compression function
  behaves as a pseudorandom family, a notion we take up in
  \cref{sec:frameworks}.

  \item[Proof of work.] Bitcoin asks miners to find a nonce making
  $\SHA(\SHA(\text{block header}))$ land below a target
  threshold~\citep{Nakamoto2008}. If the function behaves randomly, there
  is no better strategy than exhaustive guessing, so a solution \emph{is}
  verifiable evidence of expended computation. As of this writing the
  Bitcoin network evaluates the function roughly $10^{21}$ times per
  second---surely the most intensively exercised nontrivial function in
  the history of mathematics---and this entire expenditure is, among other
  things, a brute-force search for regularities that has found none.

  \item[Key derivation and pseudorandom generation.] Operating systems
  stretch and mix entropy pools through \SHA{}; standardized deterministic
  random-bit generators iterate it. Here the function is used precisely
  \emph{as} a pseudorandom number generator, which frames the question of
  \cref{sec:frameworks}: where does it sit in the space of PRNGs?
\end{description}

Note how different the demands of these uses are: fingerprinting needs
\emph{collision resistance} (no collision $x \ne y$ with $H(x) = H(y)$
has ever been \emph{published}---though note the epistemic fine print:
were one discovered, its discoverer might have excellent reasons to keep
it private),
commitments need \emph{hiding} and \emph{binding}, proof-of-work needs
\emph{progress-freeness} (no shortcut to partial preimages), key
derivation needs \emph{pseudorandomness of outputs}. A taxonomy of these
properties and their logical relations exists~\citep{RogawayShrimpton2004},
and the standard reference surveying the whole subject at textbook depth
is the Handbook chapter of \citet[Ch.~9]{MenezesOorschotVanstone1996};
the striking fact is that one fixed function is trusted, simultaneously,
with all of them.

\subsection{The original hash functions}\label{sec:elderhash}

The word ``hash'' is decades older than its cryptographic sense, and the
two meanings are worth separating before the cryptographic one takes over
the paper. Hashing as a storage technique is nearly as old as
stored-program computing: scatter-storage schemes circulated inside IBM
from 1953 (in an internal memorandum of H.~P. Luhn), reached the open
literature with Peterson's study of random-access
addressing~\citep{Peterson1957}, and were codified, history and all, in
Knuth's chapter on searching~\citep[\S 6.4]{Knuth1998}. In that tradition
a hash function need only \emph{scatter}; collisions are a fact of life
to be chained through, probed around, and amortized away, and nobody is
trying to \emph{manufacture} one.

Here is a complete, honest member of that tradition, as a working
one-liner,
\begin{center}
  \texttt{def h(k): return k \% 997}
\end{center}
---division hashing, filing the integer key $k$ into one of $997$
buckets~\citep[\S 6.4]{Knuth1998}. Mathematically it is the reduction
map $\mathbb{Z} \to \mathbb{Z}/997\mathbb{Z}$, and because that map is a
ring homomorphism, every anatomical question one could ask of it has a
one-line proof. Its collisions are not merely abundant but
\emph{listable}: the keys colliding with $k$ are exactly the arithmetic
progression $k + 997\mathbb{Z}$. It is provably injective on a natural
piece of its domain: any $997$ consecutive keys land in $997$ distinct
buckets. It is provably surjective, with exact equidistribution: among
any $997m$ consecutive keys, every bucket is hit exactly $m$ times. And
it respects algebraic structure outright:
$h(j+k) \equiv h(j) + h(k)$ and $h(jk) \equiv h(j)\,h(k) \pmod{997}$.

How far does this transparency stretch toward \SHA{}'s actual machinery?
Further than one might guess. Trade the prime modulus for a power of two
and make the map affine,
\[
  h(k) = (a\,k + b) \bmod 2^{n}, \qquad a \text{ odd},
\]
and \emph{carries} enter the story: the multiplication and addition now
ripple information from low-order bits into high ones---the $2$-adically
continuous operations we meet again in \cref{sec:statespace}. Yet nothing
is lost to analysis. With $a$ odd the map is a \emph{bijection} of the
$n$-bit words (no longer a compressing bucket hash but a reversible
mixing map), invertible in closed form as $k = a^{-1}(h - b) \bmod 2^{n}$,
where $a^{-1}$ is the inverse of $a$ modulo $2^{n}$; iterated, it is
precisely a linear congruential generator, whose period is a classical
theorem (\cref{sec:prngs}). Now throw an exclusive-or on top,
\[
  h(k) = \bigl((a\,k + b) \bmod 2^{n}\bigr) \shaxor c,
\]
and---the pleasant surprise---it stays analyzable: an $\shaxor$ with a
constant is itself a bijection, living in the \emph{rival}
$\Ftwo$-geometry of \cref{sec:statespace}, and a composition of bijections
is a bijection, so the whole map is a reversible permutation one can still
invert by hand. What one now holds is a one-round baby \SHA{}: two of its
three ARX primitives (carry-bearing arithmetic and $\shaxor$) already in
place, and analyzable only because nothing yet forces the two geometries
to \emph{interfere}. \SHA{} adds the third primitive---bitwise rotation,
the operation that couples the geometries---and iterates the mixture
$64$ times; it is that composition, not any single ingredient, that
carries the design out of reach of every one-line argument above.

\subsection{Cryptographic hash functions}\label{sec:falling}

\SHA{} is the survivor of an explicit evolutionary lineage, and the
extinctions in that lineage are the best available evidence about what
kind of object it is. The cryptographic branch split from the storage
tradition of \cref{sec:elderhash} the moment someone set out to
\emph{manufacture} collisions rather than merely tolerate them---the
moment, that is, an adversary entered the room. One-way functions
appeared first as a systems trick for protecting stored login passwords
(in use by the late 1960s; first carefully engineered constructions
in~\citealp{Purdy1974}; the classic Unix case history
in~\citealp{MorrisThompson1979}); one-wayness was then named as a
foundational primitive of the new public-key
cryptography~\citep{DiffieHellman1976}; and the modern object---a
compressing, collision-resistant, \emph{one-way hash function}---was
isolated in Merkle's 1979 dissertation~\citep{Merkle1979}. Same word as
before, escalating demand: from ``collisions are a performance nuisance''
to ``collisions exist in abundance, but no adversary on Earth can exhibit
one.'' The escalation is a change of adversary model, not of
machinery---and it took thirty years to complete.

Put the four-question anatomy exam of \cref{sec:elderhash} in front of
\SHA{} now, and every clean answer becomes an open problem.
Collisions: they exist in astronomical abundance (pigeonhole), yet none
has ever been exhibited, and no known method materially beats generic
birthday search at $2^{128}$ evaluations (\cref{sec:telescope}).
Injectivity on a restricted domain: is \SHA{} injective on the
$2^{128}$ inputs of exactly $128$ bits? For a random function this
would be very nearly a fair coin flip---the expected number of
colliding pairs is about $\tfrac12$---and for \SHA{} nobody can improve
on that shrug in either direction. Surjectivity: unknown; a random
function of these dimensions would miss at most a fantastically small
fraction of its range, but nothing of the kind is proved. Respected
structure: the design \emph{intent} is the wholesale negation of the
homomorphism property, and ``respects nothing'' is precisely the
conjunction of unproved statements this paper keeps circling. Same
species as the one-liner, same questions; but every answer has moved
from a line of algebra to an open problem. The rest of the paper lives
in that gap.

The modern era begins with Merkle's and Damg{\aa}rd's 1989 work putting
hash functions on a design principle. Rather than design a function on
arbitrary-length inputs directly, one fixes a single \emph{compression
function} $f$ that maps a fixed-size state and one message block to a new
state, and extends it to messages of any length by iterating: split the
(padded) message into blocks $m_1, \dots, m_k$, start from a fixed
initial value, and fold the blocks through $f$ one at a time,
\[
  h_0 = \mathrm{IV}, \qquad
  h_i = f(h_{i-1}, m_i), \qquad
  H(m) = h_k,
\]
reading off the final chaining value $h_k$ as the digest. The virtue of
this construction is a \emph{theorem}: any collision in the full hash $H$
can be unwound into a collision in the single map $f$, so collision
resistance of the whole reduces to collision resistance of $f$ alone (the
Merkle--Damg{\aa}rd theorem, \cref{thm:md}). This is exactly the skeleton
\SHA{} is built on, compression function and all---the finite map $f$ and
the two theorems that wrap it are anatomized in
\cref{sec:merkledamgard,sec:constructions}. Rivest's concrete proposals
MD4 (1990) and MD5 (1992)~\citep{RFC1321} were the first mass-market
instances---$128$-bit hashes built as fast iterated compression---and MD5
was, for a decade, everywhere. NIST's SHA-1 (1995) enlarged the state to
$160$ bits and adjusted the internals; it inherited MD5's ubiquity. The
SHA-2 family---of which \SHA{} is the $256$-bit member---was standardized
in 2002 with longer digests and a reworked, more conservative round
function.

Then the towers fell, in order of age. In 2005 Wang and Yu exhibited
explicit collisions in MD5 by hand-crafted differential
analysis~\citep{WangYu2005}; within a few years refinements produced
colliding documents in seconds on a laptop, and forged certificates in the
wild. The same summer, Wang, Yin, and Yu broke SHA-1's collision
resistance on paper, at cost $2^{69}$ evaluations versus the $2^{80}$
birthday bound~\citep{WangYinYu2005}; twelve years later the attack was
industrialized: the SHAttered project published two distinct PDF files
with equal SHA-1 hashes, at a cost of roughly $2^{63}$ evaluations---about
6{,}500 CPU-years plus 100 GPU-years, real but affordable for a
well-resourced organization~\citep{Stevens2017}. Git, which had bet its
object model on SHA-1, began a still-ongoing migration to \SHA{}.

Two lessons, both load-bearing for this paper. First, \emph{the choice of
algorithm matters enormously}: MD5, SHA-1, and \SHA{} look like
siblings---the same iterated-compression skeleton, the same diet of
$32$-bit additions, rotations, and Boolean operations---yet differential
cryptanalysis (anatomized in \cref{sec:differential}) shattered the first
two while \SHA{}, attacked by the same schools with the same tools for
twenty years, has not visibly moved
(current state of the art in \cref{sec:frameworks,sec:siege}). Whatever separates
broken from unbroken here is not the general species of design; it lives
in the fine structure of round counts, rotation offsets, and message
scheduling---quantitative differences producing a qualitative divide, and
nobody can currently compute, from the specification alone, which side of
the divide a design falls on. Second, breaks are not black swans; they are
the \emph{normal} fate of these objects, which makes \SHA{}'s quarter
century of stability a phenomenon in want of explanation rather than a
default to be assumed.

\subsection{The epistemological situation}\label{sec:epistemology}

How do we know what we know about \SHA{}? Not the way we know things about
the Josephus function. There is no theorem asserting any of the properties
in \cref{sec:uses}; as we will see in \cref{sec:frameworks}, for a fixed
unkeyed function it is delicate even to \emph{state} them. What exists
instead is a corpus of observations: a published specification, a fossil
record of broken relatives (\cref{sec:falling}), a large literature of
failed and partially successful attacks on weakened versions, and an
industrial base exercising the function at cosmological rates without
surfacing an anomaly.

This is the epistemic posture of a \emph{naturalist}, not an axiomatist:
we know \SHA{} the way nineteenth-century biology knew a species---by
morphology, habitat, behavior under stress, and comparison with extinct
relatives. The cryptographic literature has even formalized the ignorance:
Rogaway's ``human ignorance'' framework~\citep{Rogaway2006} proves
theorems of the form ``any explicit method for breaking protocol $P$
yields an explicit collision in \SHA{}''---transferring trust not from an
axiom but from the observed, decades-long failure of a very motivated
species to produce such a collision. It is Popperian science embedded
inside pure mathematics: the conjecture ``\SHA{} has no exploitable
structure'' is falsifiable, massively tested, and unproven.

The rest of this paper takes the naturalist's program seriously and asks
what a \emph{quantitative} natural history of this object would look like:
first the anatomy (\cref{sec:geometry}), then the field measurements one
can actually take, where work of Flajolet supplies exactly the right
instrument (\cref{sec:randommaps}), then the existing theoretical
frameworks and the genuine gap between them (\cref{sec:frameworks}).
